\documentclass[10pt,a4paper]{amsart}
\usepackage[margin=19mm]{geometry}
\usepackage{amsmath,amssymb,mathtools}
\usepackage[scr=rsfs,cal=euler]{mathalfa}
\usepackage{xfrac}
\usepackage[unicode,hidelinks,bookmarks=false]{hyperref}
\newcommand{\ie}{i.e.\ }
\newcommand{\cf}{cf.\ }
\newcommand{\resp}{respectively\ }
\newcommand{\Subsection}{\subsection}
\providecommand{\nobreakdash}{\nobreak\hbox{-}}
\setlength{\emergencystretch}{2em}
\multlinegap0pt
\usepackage{amssymb}
\usepackage{dsfont}
\usepackage{bm}
\usepackage{mathtools}
\usepackage{enumerate}
\usepackage{float}
\usepackage{xcolor}
\newtheorem{proposition}{Proposition}
\newcommand{\DE}[1]{{\color{red} #1}}
\title{{Double cluster swapping for spin models: \\Pfaffian relations and sharpness}}
\author{\ Diederik van Engelenburg$^*$ \and Lorca Heeney$^*$ \and Marcin Lis}
\date{Source: arXiv:2609.01362v1 (CC BY 4.0)}
\begin{document}
\maketitle

\noindent\emph{Source and licence.} {Double cluster swapping for spin models: \\Pfaffian relations and sharpness}. Version arXiv:2609.01362v1 (CC BY 4.0), distributed by arXiv under the Creative Commons Attribution 4.0 International licence, http://creativecommons.org/licenses/by/4.0/, as recorded in the arXiv licence metadata for this exact version and shown on the abstract page https://arxiv.org/abs/2609.01362v1. Full licence notice: \texttt{licenses/arxiv-2609.01362-CC-BY-4.0.txt}.\par\smallskip

\noindent\textit{Editorial scope.} Counted unit: the article's proposition prop:main (source0.tex lines 1128--1131). The article's complete original proof of it is delivered with it (source0.tex lines 1788--1892, one \texttt{proof} environment of 105 lines, which the article places away from the statement). No mathematics is written, repaired or re-derived: the statement and the proof are the article's own lines, in the article's own notation, and no retained line is substituted or re-flowed.  No further source-local context is needed: every cross-reference of every retained line resolves inside the two delivered spans.  Nothing was omitted from any retained span, so the package carries no omission record.  No retained line contains a citation, so the wrapper carries no bibliography and no bibliography entry.  No retained line contains a figure environment, an image-inclusion command or drawing code, so no image is delivered and the package declares no asset dependency.  Editorial formatting only, in addition: an AMS \texttt{amsart} preamble that restates the article's own declarations and macros used by the retained lines, namely the environments \texttt{proposition} on separate counters, together with the standard packages those lines need.  The retained environments print in this document as Proposition 1 in order of appearance, whereas the article prints the same items with the numbering of its own full document.  The complete native source of the article as distributed in the arXiv e-print for this exact version is bundled unchanged as \texttt{sources/209103/source0.tex}.
\begin{proposition} \label{prop:main}
    Let $G = (V, E)$ be a finite graph with a distinguished tuple $W$ of vertices. 
    If $G$ satisfies the $W$-crossing property and is fully reduced, then $G$ can be drawn in the disk. 
\end{proposition}

\begin{proof}[Proof of Proposition \ref{prop:main}]
    We will apply induction on the number of vertices of $G = (V, E)$. 
    For the base case, $|V| = |W| = 4$, the result is trivial. 
    %\DE{Maybe it would be a bit easier to do the base case also for one internal vertex. In that case, we can assume two internal vertices and the contracted thing would be completely contained in $V \setminus W$. As such, the fact that $G / e$ being drawn in the disk is really enough, but maybe it is not necessary to separate it. }
    
    Assume the result holds for all finite graphs with $ k - 1$ or less vertices.
    Let $G = (V, E)$ be a finite graph with $k$ vertices, which is fully reduced. 
    If $|V| = |W|$, the result is an immediate consequence of the $W$-crossing property and we are done. 
    Thus, we may suppose without loss of generality that there is some vertex $x \in V \setminus W$. 
    Pick any such vertex and an edge $e = \{x, y\}$ emanating from it.
    
    Write $G / e$ for the graph obtained from contracting $e $, let $v_e$ be the obtained vertex. 
    Then $G / e$ has $k - 1$ vertices, hence the induction hypothesis implies that either $G / e$ can be reduced, 
    or that $G / e$ can be drawn in the disk. We handle the two cases separately. 
	\\
	\\
	\noindent \textbf{Case 1}. 
	$G / e$ can be reduced, with reduction set $S$.
    Write $X$ for the (non-empty) maximal disconnected set. 
    
    Notice that $v_e \in S$, as otherwise $S$ would be a reduction set for $G$, but $G$ is fully reduced by assumption. 
    A similar argument gives $|S| = 3$. 
    Hence, in the graph $G$, the set $P = (S \setminus \{v_e\}) \cup \{x, y\}$ has size four and $P$ is a $(W, X)$ separator. 
    
    Note that $P$ may contain elements of $W$, but by construction is not equal to $W$. 
    Let $Y$ be the connected component of $W$ in $G - P$, so $V$ is the disjoint union of $X, P$ and $Y$. 
    Thus, by construction, both $X$ and $Y$ are non-empty. 
    
    Write $G_X$ for the subgraph of $G$ induced by $X \cup P$ and $G_Y$ for the one induced by $Y \cup P$. 
    
    By Menger's theorem, there exist $4$ disjoint paths connecting $W$ to $P$ in $G - P$, 
    since $G$ is fully reduced. 
    As such, the ordering of elements in $W$ induces and ordering of $P = (p_1, p_2, p_3, p_4)$, 
    and the graph $G_X$ must satisfy the $P$-crossing property. 
    Moreover, $G_X$ is reduced because $G$ is reduced by assumption. 
    Since $|X \cup P| < |V|$, the induction hypothesis implies that $G_X$ can be drawn in the disk with $P$ on the boundary of the disk. 
    
    Write next $G_Y'$ for the graph obtained by adding to $G_Y$ edges between elements of $P$ in cyclic order. 
    We make the following claim. 
    \\
    \\
    \noindent \textbf{Claim:}
    The graph $G_Y'$ is fully reduced and satisfies the $W$-crossing property. 
    \\
    \\
    \noindent Let us first argue how to conclude the proof of Proposition \ref{prop:main} from this claim. 
    Since the number of vertices of $G_Y'$ is less or equal to $k -1$, 
    the claim and the induction hypothesis imply that $G_Y'$ can be drawn in the disk. 
    The vertices of $P$ lie on one cycle, hence they must form a face in the drawing. 
    Indeed, consider four disjoint paths ${\gamma}_1, \ldots, {\gamma}_4$ connecting 
    $x_1, \ldots, x_4$ to $p_1, \ldots, p_4$ inside $G_Y'$ 
    % (which must exist, and do not use the edges connecting $p_1, \ldots, p_4$ that were added to create $G_Y'$. 
    If $P$ is not a face, then we must have (up to relabeling) 
    that the edge connecting $p_1, p_2$ either intersects one of the paths ${\gamma}_1, \ldots, {\gamma}_4$, 
    intersects one of the other edges $p_i, p_{i+1}$ or must go outside of the disk 
    (since $W$ lies on the boundary of the disk). 
    None of this can happen because $G_Y'$ is drawn in the disk with $W$ on the boundary, so $P$ is a face. 
    We can put $G_X$ inside this graph by gluing $G_X$ and $G_Y$ at the vertices of $P$ 
    and the obtained graph is planar and can be drawn in the disk. 
    Finally, removing the edges in the cycle of $P$ from this graph, we get back $G$ and hence for Case 1, 
	the induction is finished.
    
    We are left to prove the claim. 
	Observe that $G_Y'$ is fully reduced as any reduction set would also be a reduction set of $G$. 
    % We first show that $G_Y'$ is fully reduced. 
    % Suppose not, then there would be some non-empty set $X'$ of vertices such that there is 
    % a $(W, X)$ separator of size at most $3$. 
    % But then this is also a separator of $G_Y$ and hence of $G$, 
    % which contradicts the assumption that $G$ cannot be reduced.
    
    For the $W$ crossing property, argue by contradiction. 
	Suppose there exist two simple, disjoint paths in $G_Y'$ violating the $W$-crossing property, 
	say $\gamma$ and $\widetilde{\gamma}$ and say they connect $x_1, x_3$ and $x_2, x_4$ in $W$ respectively. 
	Both $\gamma$ and $\widetilde{\gamma}$ use (at least) one edge added between the vertices of $P$. 
	Indeed, suppose that $\gamma$ does not use any of the edges added to $P$, 
	then it is a path in $G$. 
	Since $G_X$ is connected, also the path $\widetilde{\gamma}$ can be lifted to $G$ without intersecting $\gamma$, 
	contradicting the $W$-crossing property of $G$. 

	Thus, $\gamma$ and $\widetilde{\gamma}$ use an edge of $P$ and all elements of $P$, 
	and since $G_Y$ is planar and can be drawn in the disk with $W$ as a single face, the induced ordering on $P$ by $\gamma, \widetilde{\gamma}$
	viewed in $G_Y$ has to agree with the one induced by $W$ above. 
	This is a contradiction with the existence of the non-intersecting simple paths $\gamma$ and $\widetilde{\gamma}$,
	and as such proves the claim. 
	\\
	\\
	\textbf{Case 2:} $G / e$ is fully reduced. 
	By the induction hypothesis, $G / e$ can be drawn in the disk with $W$ on the boundary. 
	The only problem that can still occur is that when `unmerging' $xy$, there are edges that now must cross. 
	Consider the set $N_x$ of neighbors of $x$ in $G$. 
	The graph $G - x$ is planar since $G / e$ is planar, and thus can be drawn with $W$ on the boundary of the disk. 
	We claim that $N_x$ must lie on a face of $G - x$, in which case we are done. 
	Suppose not. 
	For every subset $A \subset N_x$ of size $3$, there must exist $3$ disjoint paths simple paths to $W$ in $G - x$, 
	since $G$ is fully reduced, and furthermore, there exists a further simple disjoint path from $x$ to $W$, 
	which exits $N_x$ at $a'$. 
	This induced an ordering of the elements in $A \cup \{a'\} = (a_1, \ldots, a_4)$. 
	
	If $N_x$ does not lie on a face of $G - x$, there hence exists an ordered set $A' = (a_1, \ldots, a_4)$ contained in $N_x$, 
	such that there are four disjoint paths $\tau_1, \ldots, \tau_4$ connecting $A'$ to $W$ \emph{and furthermore} such that
	(up to cyclic relabeling) there exists a further path in $G - x$ connecting $a_1$ and $a_3$ which only intersects $\tau_1, \ldots, \tau_4$ 
	at $a_1, a_3$. 
	But this readily implies that in $G$, the $W$-crossing property is violated, so $N_x$ must lie on a face of $G - x$,
	finishing the induction in Case 2 and therefore the proof of Proposition \ref{prop:main}. 
\end{proof}

\end{document}
